\justifying
\NSFCBodyText

\subsubsection{与本项目相关的研究工作积累}

申请人长期从事多模态信息处理、视觉问答与跨模态检索，并与新疆大学民族语言信息处理方向合作开展图像理解研究。与本课题直接相关的积累有三类。第一，知识增强与空间关系建模的视觉问答，能够把问题条件作用于多层视觉表征\cite{yan2024knowledge,yan2022spca}。第二，遥感视觉问答中的视觉校准与状态空间推理，说明申请人具备在专业图像上训练多模态理解模型的经验\cite{xu2026rshr,xu2026rssr}。第三，跨模态检索中的特征融合，为图文对齐与检索评价提供方法基础。新能源功率预测与小目标检测等工作\cite{sun2025dwt,liu2026wavekan,wang2025syolo}证明申请人具备深度学习训练与实验组织能力，但不构成本课题的研究对象。

\subsubsection{已取得的研究工作成绩}

申请人已以第一或通讯作者发表多篇 SCI 收录论文，满足实验室对申请人的资格要求。代表性工作包括 Machine Learning、The Visual Computer、Expert Systems with Applications 与 Computer Vision and Image Understanding 等期刊上的视觉问答与图像理解论文\cite{yan2024knowledge,yan2022spca,xu2026rshr,xu2026rssr}。上述成绩用于证明有能力完成本课题的编码器适配、融合训练与检索评价，不把其中的新能源或遥感指标表述为本课题已经完成的结果。

\subsubsection{项目组能力与任务分工}

申请人颜丰负责总体设计、方法实现与论文撰写。课题组现有约 10 名研究生，可承担数据整理、基线复现、消融实验与评测记录。任务按研究内容划分：场景文字编码器与 RUST 测试、融合与对齐训练、跨语言检索与理解题核对。不把未书面确认的实验室固定人员写成课题组成员。

\subsubsection{前期基础与本项目的衔接}

内容一衔接申请人在视觉表征与检测方面的训练经验，用于场景文字区域编码。内容二衔接视觉问答中的跨模态交互与检索中的特征融合。内容三衔接已有的检索/问答评价习惯，改为 Recall@$K$、MRR 与小规模理解准确率。天池英才项目研究产业集群场景下的多模态融合，与本课题的民族语言场景图文对象不同，不构成重复申报。
